% answers_short.tex —— 2008 年·普通高等学校招生考试·江西卷·文科 · 简短答案速查
% 由 tools/make_short_answers.py 自 answers.tex 的 \gksolhead/\gkanswer 提取，
% 并与 2008/2008_数学真题_江西文_答案要点.md 逐题比对后生成，勿手改（改源 MD 后重跑）。

\begin{gknotice}
\textbf{使用说明}：本节只给最终答案便于快速核对，完整推导见解析版。答案为编写者独立推导，\textbf{非官方标准答案}。
\end{gknotice}

\gkpart{选择题}
\begin{center}\begin{tabular}{ccccccccccccc}
\toprule
题号 & 1 & 2 & 3 & 4 & 5 & 6 & 7 & 8 & 9 & 10 & 11 & 12 \\
\midrule
答案 & B & D & B & C & A & A & C & D & B & D & C & C \\
\bottomrule
\end{tabular}\end{center}

\gkpart{填空题}
\begin{itemize}[leftmargin=2.6em]
  \item[\textbf{13.}] $[-3,1]$。
  \item[\textbf{14.}] $\dfrac{x^2}4-\dfrac{3y^2}4=1$。
  \item[\textbf{15.}] $5$。
  \item[\textbf{16.}] ABD。
\end{itemize}

\gkpart{解答题}
\begin{itemize}[leftmargin=2.6em]
  \item[\textbf{17.}] （1）$1$；（2）$\sqrt5$。
  \item[\textbf{18.}] （1）$0.20$；（2）$0.48$。
  \item[\textbf{19.}] （1）$a_n=2n+1$，$b_n=8^{n-1}$；（2）$\dfrac{n(3n+5)}{4(n+1)(n+2)}$。
  \item[\textbf{20.}] （1）证明见解析；（2）$\arccos\dfrac{\sqrt6}6$。
  \item[\textbf{21.}] （1）增区间 $(-2a,0)$、$(a,+\infty)$，减区间 $(-\infty,-2a)$、$(0,a)$；（2）$(0,1)\cup\left(\sqrt[4]{\dfrac{12}7},+\infty\right)$。
  \item[\textbf{22.}] （1）证明见解析；（2）存在，$y_0>1$，距离为 $1$。
\end{itemize}
